\section{Experiments}
\label{sec:experiments}

\subsection{Experimental Setup}
\label{subsec:setup}

We conduct two complementary experiments to evaluate IC-Knock-Poly under different data regimes:

\textbf{Experiment 1: Independent Base Features.} Base features are independent with label noise $\sigma = 0.5$.

\textbf{Experiment 2: Correlated Base Features with Label Noises.} Base features have block-diagonal correlation ($\rho=0.8$) with high label noise ($\sigma = 3.0$).

\textbf{Common Setup.} Each dataset has $n$ training samples, $n_{\text{val}} = 100$ validation samples, and $n_{\text{test}} = 200$ test samples. We vary: $n \in \{50, 75, 100\}$, $p \in \{1, 3, 5, 7, 9\}$, $k \in \{2, 3, 5\}$, and polynomial degree $d \in \{2, 3\}$.

\subsection{Main Results}
\label{subsec:main-results}

Table~\ref{tab:exp1-results} and Table~\ref{tab:exp2-results} compare validation-based methods on two configurations: moderate dimensionality ($p=5$) and high dimensionality ($p=9$), with fixed $n=100, k=3, d=3$.

\begin{table}[htbp]
\centering
\caption{Experiment 1: Independent Base Features ($n=100, k=3, d=3$, noise=0.5, 3 trials). Values show mean$\pm$std for FDR/TPR, and mean for $R^2$/RMSE/\#Selected.}
\label{tab:exp1-results}
\small
\begin{tabular}{lccccc}
\toprule
\textbf{Method} & \textbf{FDR} & \textbf{TPR} & \textbf{$R^2$} & \textbf{RMSE} & \textbf{\#Sel} \\
\midrule
\multicolumn{6}{l}{\textbf{Moderate Dimensionality} ($p=5$)} \\
\midrule
IC-Knock-Poly-Val & \textbf{0.48$\pm$0.42} & \textbf{0.78$\pm$0.39} & 0.650 & 0.869 & 7.7 \\
Poly-Knockoff-Val & 0.85$\pm$0.13 & 0.22$\pm$0.19 & 0.649 & 0.872 & 7.0 \\
Poly-Lasso-Val & 0.92$\pm$0.08 & 0.22$\pm$0.19 & 0.660 & 1.229 & 8.0 \\
Poly-CLIME-Val & 0.81$\pm$0.17 & 0.22$\pm$0.19 & 0.656 & 1.383 & 6.0 \\
Poly-OMP-Val & 0.83$\pm$0.29 & 0.11$\pm$0.19 & 0.654 & 1.682 & 7.0 \\
\midrule
\multicolumn{6}{l}{\textbf{High Dimensionality} ($p=9$)} \\
\midrule
IC-Knock-Poly-Val & 0.79$\pm$0.07 & \textbf{0.78$\pm$0.19} & \textbf{0.980} & 1.073 & 14.7 \\
Poly-Knockoff-Val & 0.81$\pm$0.17 & 0.33$\pm$0.33 & 0.978 & 1.203 & 11.0 \\
Poly-Lasso-Val & 0.93$\pm$0.09 & 0.33$\pm$0.33 & 0.978 & 1.365 & 21.0 \\
Poly-CLIME-Val & 0.81$\pm$0.17 & 0.33$\pm$0.33 & 0.976 & 1.801 & 10.3 \\
Poly-OMP-Val & 0.61$\pm$0.35 & 0.33$\pm$0.33 & 0.981 & 2.000 & 3.0 \\
\bottomrule
\end{tabular}
\end{table}

\begin{table}[htbp]
\centering
\caption{Experiment 2: Correlated Features with Label Noises ($n=100, k=3, d=3$, noise=3.0, 3 trials). Values show mean$\pm$std for FDR/TPR, and mean for $R^2$/RMSE/\#Selected.}
\label{tab:exp2-results}
\small
\begin{tabular}{lccccc}
\toprule
\textbf{Method} & \textbf{FDR} & \textbf{TPR} & \textbf{$R^2$} & \textbf{RMSE} & \textbf{\#Sel} \\
\midrule
\multicolumn{6}{l}{\textbf{Moderate Dimensionality} ($p=5$)} \\
\midrule
IC-Knock-Poly-Val & 0.83$\pm$0.00 & \textbf{0.44$\pm$0.19} & \textbf{0.839} & \textbf{3.68} & 8.0 \\
Poly-Knockoff-Val & 0.92$\pm$0.14 & 0.11$\pm$0.19 & 0.694 & 7.45 & 6.0 \\
Poly-Lasso-Val & 0.97$\pm$0.05 & 0.11$\pm$0.19 & 0.527 & 6.30 & 15.0 \\
Poly-CLIME-Val & 1.00$\pm$0.00 & 0.00$\pm$0.00 & 0.683 & 6.89 & 5.3 \\
Poly-OMP-Val & 0.83$\pm$0.29 & 0.11$\pm$0.19 & 0.689 & 8.29 & 2.0 \\
Poly-STLSQ-Val & 1.00$\pm$0.00 & 0.00$\pm$0.00 & -32.6 & 53.1 & 36.0 \\
\midrule
\multicolumn{6}{l}{\textbf{High Dimensionality} ($p=9$)} \\
\midrule
IC-Knock-Poly-Val & 0.74$\pm$0.02 & \textbf{0.56$\pm$0.19} & \textbf{0.996} & \textbf{3.08} & 17.7 \\
Poly-Knockoff-Val & 0.83$\pm$0.29 & 0.11$\pm$0.19 & 0.995 & 4.54 & 3.3 \\
Poly-Lasso-Val & 0.99$\pm$0.01 & 0.11$\pm$0.19 & 0.932 & 13.50 & 55.0 \\
Poly-CLIME-Val & 0.89$\pm$0.19 & 0.11$\pm$0.19 & 0.995 & 4.38 & 4.0 \\
Poly-OMP-Val & 0.89$\pm$0.19 & 0.11$\pm$0.19 & 0.962 & 26.21 & 3.7 \\
Poly-STLSQ-Val & 0.98$\pm$0.04 & 0.11$\pm$0.19 & -1.34 & 79.37 & 174.3 \\
\bottomrule
\end{tabular}
\\[3pt]
\footnotesize{Note: FDR/TPR show mean$\pm$std. Best values in bold.}
\end{table}

\textbf{Key Findings:}
\begin{enumerate}
    \item \textbf{Variable Selection:} IC-Knock-Poly-Val achieves 3.5$\times$ higher TPR than baselines in Exp 1 (0.78 vs 0.22 at $p=5$), and 4$\times$ higher in Exp 2 (0.44 vs 0.11 at $p=5$).
    \item \textbf{High-Dimensional Advantage:} In Exp 2, increasing $p$ from 5 to 9 maintains IC-Knock-Poly-Val's TPR (0.44$\rightarrow$0.56) while improving $R^2$ (0.839$\rightarrow$0.996) and reducing RMSE (3.68$\rightarrow$3.08).
    \item \textbf{Prediction Quality:} IC-Knock-Poly-Val achieves substantially better prediction than Lasso in Exp 2 ($R^2$: 0.839 vs 0.527 at $p=5$; 0.996 vs 0.932 at $p=9$).
    \item \textbf{Model Sparsity:} IC-Knock-Poly-Val selects 8--18 features vs Lasso's 15--55 in Exp 2, avoiding overfitting.
\end{enumerate}

\subsection{Results under Extreme Conditions}
\label{subsec:extreme}

To evaluate robustness under severely degraded data quality, we conduct an extreme experiment with challenging noise structures: non-linear heteroscedasticity (variance $\propto \|\mathbf{x}\|^{1.5}$), 10\% outlier contamination (Student-$t$, df=2), and systematic feature-dependent bias. Data features block-diagonal correlation with $\rho=0.99$ intra-block correlation, approaching multicollinearity.

\begin{table}[htbp]
\centering
\caption{Extreme conditions: n=50, p=9 (small sample, high dimension). Mean$\pm$std over 3 trials.}
\label{tab:extreme-high-dim}
\small
\begin{tabular}{lccccc}
\toprule
\textbf{Method} & \textbf{Test $R^2$} & \textbf{TPR} & \textbf{FDR} & \textbf{\#Selected} & \textbf{Best Param} \\
\midrule
IC-Knock-Poly-Val & \textbf{0.25$\pm$0.06} & \textbf{0.13$\pm$0.19} & 0.89$\pm$0.16 & 4.3$\pm$2.5 & Q=0.05 \\
Poly-Knockoff-Val & -0.36$\pm$1.03 & 0.00$\pm$0.00 & 1.00$\pm$0.00 & 4.0$\pm$2.6 & Q=0.10 \\
Poly-Lasso-Tuned & -3.23$\pm$6.09 & 0.00$\pm$0.00 & 1.00$\pm$0.00 & 5.7$\pm$2.5 & $\alpha$=1.0--10.0 \\
Poly-CLIME-Tuned & -15.68$\pm$21.5 & 0.00$\pm$0.00 & 0.67$\pm$0.47 & 1.3$\pm$1.5 & $\alpha$=1.0--2.0 \\
Poly-OMP-Tuned & \multicolumn{5}{c}{\textit{All hyperparameters failed}} \\
Poly-STLSQ-Tuned & \multicolumn{5}{c}{\textit{All thresholds failed}} \\
\bottomrule
\end{tabular}
\end{table}

\begin{table}[htbp]
\centering
\caption{Extreme conditions: n=100, p=5 (moderate dimension). Mean$\pm$std over 3 trials.}
\label{tab:extreme-moderate}
\small
\begin{tabular}{lccccc}
\toprule
\textbf{Method} & \textbf{Test $R^2$} & \textbf{TPR} & \textbf{FDR} & \textbf{\#Selected} & \textbf{Best Param} \\
\midrule
IC-Knock-Poly-Val & -0.21$\pm$0.31 & \textbf{0.13$\pm$0.09} & 0.72$\pm$0.21 & 2.0$\pm$1.0 & Q=0.05 \\
Poly-Knockoff-Val & -0.01$\pm$0.08 & 0.00$\pm$0.00 & 1.00$\pm$0.00 & 1.7$\pm$0.6 & Q=0.10 \\
Poly-Lasso-Tuned & 0.01$\pm$0.38 & 0.00$\pm$0.00 & 1.00$\pm$0.00 & 4.0$\pm$1.7 & $\alpha$=1.0 \\
Poly-CLIME-Tuned & -1.50$\pm$1.42 & 0.00$\pm$0.00 & 0.67$\pm$0.47 & 1.3$\pm$1.2 & $\alpha$=0.5--5.0 \\
Poly-OMP-Tuned & \multicolumn{5}{c}{\textit{All hyperparameters failed}} \\
Poly-STLSQ-Tuned & \multicolumn{5}{c}{\textit{All thresholds failed}} \\
\bottomrule
\end{tabular}
\end{table}

\textbf{Key Findings under Extreme Conditions:}
\begin{enumerate}
    \item \textbf{OMP and STLSQ collapse completely:} All hyperparameter values fail to produce valid models, indicating these methods are fundamentally incompatible with extreme noise structures.
    \item \textbf{IC-Knock-Poly-Val is the only method achieving TPR$>$0:} While TPR remains low (0.13), it is the sole method identifying any true polynomial terms. Baselines universally achieve TPR=0 even with extensive hyperparameter tuning.
    \item \textbf{FDR patterns reveal failure modes:} CLIME exhibits bimodal FDR behavior---either selecting nothing (FDR=0) or selecting exclusively false features (FDR=1.0)---averaging to 0.67. This indicates the method cannot reliably distinguish signal from noise under extreme conditions.
    \item \textbf{Prediction quality:} IC-Knock-Poly-Val maintains positive $R^2$ (0.25) in the high-dimensional setting (n=50, p=9), while all baselines achieve negative $R^2$, performing worse than predicting the mean.
\end{enumerate}

\paragraph{Why Methods Fail Differently.}
\begin{itemize}
    \item \textbf{OMP and STLSQ:} Orthogonal matching pursuit and sequential thresholding require stable correlation structures. Extreme heteroscedasticity and outliers break their greedy selection criteria, causing numerical instability across all hyperparameter settings.
    \item \textbf{Lasso with tuning:} While hyperparameter tuning improves $R^2$ from $-26$ (fixed $\alpha$) to $0.01$ (tuned), TPR remains 0. The method finds predictive linear combinations but cannot identify the true sparse polynomial structure.
    \item \textbf{CLIME:} The precision-matrix-based approach fails when $\rho=0.99$ creates near-singular covariance matrices. CLIME exhibits bimodal failure behavior: in 2 of 3 trials it selects nothing (empty selection, FDR=0 by convention), while in the remaining trial it selects 3 features that are all false positives (FDR=1.0). This produces the counter-intuitive combination of TPR=0 (no true features ever selected) with FDR=0.67 (average of 0 and 1.0). The method cannot reliably distinguish signal from noise under extreme conditions.
    \item \textbf{IC-Knock-Poly-Val:} GMM-based knockoff generation provides distributional robustness against heavy-tailed noise, while conditional knockoffs adapt to local correlation structures. The validation mechanism automatically selects conservative Q values (Q=0.05 in all trials) appropriate for extreme noise.
\end{itemize}

\subsection{High-Dimensional Analysis: The Dimension Benefit}
\label{subsec:high-dim}

Table~\ref{tab:varying-p} shows IC-Knock-Poly-Val performance across feature dimensions in a challenging setting: small sample size ($n=50$), high noise ($\sigma=3.0$), and correlated features ($\rho=0.8$).

\begin{table}[htbp]
\centering
\caption{Performance vs. feature dimension $p$ on fixed configuration ($n=50, k=3, d=3$, noise=3.0).}
\label{tab:varying-p}
\small
\begin{tabular}{lcccc}
\toprule
$p$ & \textbf{FDR} & \textbf{TPR} & \textbf{Test $R^2$} & \textbf{\#Selected} \\
\midrule
1 & 0.00 & 0.56 & 0.23 & 1.7 \\
3 & 0.52 & 0.33 & -0.01 & 4.7 \\
5 & 0.83 & 0.44 & 0.84 & 8.0 \\
7 & 0.66 & 0.56 & 0.99 & 6.3 \\
9 & 0.74 & 0.56 & 1.00 & 17.7 \\
\bottomrule
\end{tabular}
\end{table}

\textbf{Counter-Intuitive Finding:} IC-Knock-Poly-Val's prediction accuracy \textit{improves} with feature dimension: $R^2 = 0.23$ at $p=1$ (failure), but $R^2 \geq 0.84$ when $p \geq 5$, reaching $R^2 = 1.00$ at $p=9$.

\paragraph{Why Dimension Helps.}
\begin{enumerate}
    \item \textbf{Ensemble Effect:} Block-diagonal correlation ($\rho=0.8$) provides redundant information, acting as implicit ensemble averaging against noise.
    \item \textbf{Improved GMM Estimation:} At $p \leq 3$, insufficient data prevents accurate covariance estimation. As $p$ increases, GMM better characterizes correlation structure.
    \item \textbf{Signal Amplification:} At $p \geq 5$, accurate GMM modeling enables better distinction between signal and noise.
\end{enumerate}

\subsection{Sample Size Analysis}
\label{subsec:sample-size}

Table~\ref{tab:varying-n} examines how IC-Knock-Poly-Val benefits from additional training data ($p=5, k=3, d=3$, noise=3.0).

\begin{table}[htbp]
\centering
\caption{Performance vs. sample size $n$ on fixed configuration ($p=5, k=3, d=3$, noise=3.0).}
\label{tab:varying-n}
\small
\begin{tabular}{lcccc}
\toprule
$n$ & \textbf{FDR} & \textbf{TPR} & \textbf{Test $R^2$} & \textbf{\#Selected} \\
\midrule
50 & 0.67 & 0.33 & 0.80 & 3.7 \\
75 & 0.84 & 0.33 & 0.99 & 6.7 \\
100 & 0.83 & 0.44 & 0.84 & 8.0 \\
\bottomrule
\end{tabular}
\end{table}

Prediction accuracy improves substantially from $n=50$ to $n=75$ ($R^2$: 0.80$\rightarrow$0.99). The slight decrease at $n=100$ reflects variance from small trial counts.

\subsection{Validation-Based Selection}
\label{subsec:validation}

Table~\ref{tab:fixed-q} shows the validation process on configuration ($n=100, p=5, k=3, d=3$, noise=3.0).

\begin{table}[htbp]
\centering
\caption{Fixed-Q candidates vs. Validation-selected performance ($n=100, p=5, k=3, d=3$, noise=3.0).}
\label{tab:fixed-q}
\small
\begin{tabular}{lccc}
\toprule
\textbf{Variant} & \textbf{FDR} & \textbf{TPR} & \textbf{Test $R^2$} \\
\midrule
Q0.05 (candidate) & 0.84 & 0.33 & 0.84 \\
Q0.10 (candidate) & 0.84 & 0.33 & 0.84 \\
Q0.15 (candidate) & 0.83 & 0.44 & 0.84 \\
\midrule
\textbf{Val} (selected) & \textbf{0.83} & \textbf{0.44} & \textbf{0.84} \\
\bottomrule
\end{tabular}
\end{table}

The validation procedure correctly identifies Q0.15 as optimal (highest TPR with same $R^2$), demonstrating automatic hyperparameter selection.
